\documentclass[10pt,a4paper]{article}
\usepackage{amsmath,amssymb,amsthm,graphicx}
\title{Oracle separation of error-prone easiness and errorless average-case hardness of NP}
\author{Shuichi Hirahara, Mikito Nanashima}
\date{doi:10.4230/LIPIcs.CCC.2022.25; CC BY4.0}
\begin{document}
\maketitle
\noindent\textit{Editorial note.} Original authors: Shuichi Hirahara and Mikito Nanashima. Official CCC2022 article25 (2022-07-11), DOI 10.4230/LIPIcs.CCC.2022.25. Creative Commons Attribution4.0.\par\medskip
\noindent\textit{Editorial note.} License: https://creativecommons.org/licenses/by/4.0/\par\medskip
\noindent\textit{Editorial note.} Faithful editable TeX transcription of the two core average-case separation clauses of author Theorem4, complete masked-oracle construction, Theorem15/Lemmas16-17 with full AppendixA, Theorem19 and Corollaries21-22 with all proof continuations, necessary definitions and complete bibliography. Original author mathematical expressions retained as published. Remaining meta-complexity/learning/cryptographic consequence clauses not separately counted. Standalone typography and long-formula line breaks adapted; no new proofs or mathematical corrections.\par\medskip
\noindent\textit{Allowed standard prerequisites (editorial).}\begin{itemize}
\item Baby switching lemma and pseudorandom generation from a random function oracle. Author [14] switching lemma and exact Theorem18 citing Zimand[45] retained as standard prior results.
\item AIPRG-to-hitting-set transformation and Borel-Cantelli lemma. Author invokes [34, Lemma18] for the standard primitive conversion; every local oracle proof and errorless-hardness reduction is retained.
\end{itemize}\medskip
\paragraph{Author notation and average-case conventions (Section 3).}
For each $n\in\mathbb N$, let $[n]=\{1,\ldots,n\}$. For each $x\in\{0,1\}^n$ and $i\in[n]$, we let $x_i$ denote the $i$-th bit of $x$ and $x_{\leq i}$ denote $x_1\circ\cdots\circ x_i$. For a distribution $D$, we write $x\leftarrow D$ to refer to a random sampling $x$ according to $D$. For a finite set $S$, we also use the notation $x\leftarrow_u S$ to denote the uniform sampling from $S$. For each $n\in\mathbb N$, we let $U_n$ denote the uniform distribution over $\{0,1\}^n$ or a random variable selected uniformly at random from $\{0,1\}^n$ in context. We use the notation $\operatorname{negl}$ to represent a certain negligible function, i.e., for any polynomial $p(n)$, $\operatorname{negl}(n)<1/p(n)$ for sufficiently large $n\in\mathbb N$. For a randomized algorithm $A$ using $r(n)$ random bits on an $n$-bit input, we use $A(x;s)$ to refer to the execution of $A(x)$ with a random tape $s$ for $x\in\{0,1\}^n$ and $s\in\{0,1\}^{r(n)}$.
For any oracles $O_0$ and $O_1$, we let $O_0+O_1$ denote the combination, i.e., for any $b\in\{0,1\}$ and any $x\in\{0,1\}^*$, $(O_0+O_1)(b\circ x)=O_b(x)$.
In this paper, we assume the basic knowledge of probability theory, including the union bound, Markov's inequality, Hoeffding's inequality, and the Borel--Cantelli lemma.
For each $p\in[0,1]$ and set $S$ of variables taking binary values, we define a $p$-random restriction $\rho$ to $S$ as a partial assignment $\rho:S\to\{0,1,*\}$ (where $*$ represents ``unassigned'') randomly selected as follows: for each $x\in S$,
\[\rho(x)=\begin{cases}*&\text{with probability }p,\\0&\text{with probability }(1-p)/2,\\1&\text{with probability }(1-p)/2.\end{cases}\]
For every restriction $\rho$ to $S$ and function $f$ defined on $S$, we let $f|_\rho$ denote the restricted function obtained by applying a partial assignment to $f$ according to $\rho$.
\paragraph{3.1. Average-Case Complexity.}
We present the notions in average-case complexity theory. Further backgrounds can be found in a survey [7].
We say that a family $D=\{D_n\}_{n\in\mathbb N}$ of distributions, where each $D_n$ is a distribution on $\{0,1\}^n$, is (polynomial-time) samplable if there exists a randomized sampling algorithm $S$ such that the distribution of $S(1^n)$ is identical to $D_n$ for each $n\in\mathbb N$. We consider a family of distributions as a single distribution on instances. We define a distributional problem as a pair of a language $L\subseteq\{0,1\}^*$ and a distribution $D=\{D_n\}_{n\in\mathbb N}$ on instances. For a standard complexity class $C$ (e.g., $\mathsf{NP}$), we define its average-case extension $\mathsf{Dist}C$ as $\mathsf{Dist}C=\{(L,D):L\in C,D\text{ is samplable}\}$.
We present the errorless average-case easiness. We say that a distributional problem $(L,D)$ has an errorless heuristic algorithm $A$ with failure probability $\epsilon:\mathbb N\to(0,1)$ if (1) $A$ outputs $L(x)(:=\mathbf1\{x\in L\})$ or $\bot$ (which represents ``failure'') for every $x\in\operatorname{supp}(D)$, and (2) the failure probability that $A(x)$ outputs $\bot$ over the choice of $x\leftarrow D$ is bounded above by $\epsilon(n)$ for each $n\in\mathbb N$. Note that an errorless heuristic algorithm never outputs an incorrect value $\neg L(x)$ for any $x\in\operatorname{supp}(D)$. Then, for every $\epsilon:\mathbb N\to(0,1)$, we define a class $\mathsf{Avg}_\epsilon\mathsf P$ as a class of distributional problems that have a polynomial-time errorless heuristic algorithm with failure probability $\epsilon(n)$. Furthermore, we say that a distributional problem $(L,D)$ has an errorless heuristic scheme $A$ if $A$ is given an instance $x\in\operatorname{supp}(D)$ and $\epsilon\in(0,1)$ as input and satisfies the condition of an errorless heuristic algorithm with failure probability $\epsilon$. We define a class $\mathsf{AvgP}$ as a class of distributional problems that have a polynomial-time errorless heuristic scheme. It is not hard to verify that $\mathsf{AvgP}\subseteq\mathsf{Avg}_{1/p(n)}\mathsf P$ for any polynomial $p(n)$.
Next, we present the error-prone average-case easiness. We say that a distributional problem $(L,D)$ has an error-prone heuristic algorithm $A$ with failure probability $\epsilon:\mathbb N\to(0,1)$ if the failure probability that $A(x)\neq L(x)$ over the choice of $x\leftarrow D$ is bounded above by $\epsilon(n)$ for each $n\in\mathbb N$. Note that an error-prone heuristic algorithm may output an incorrect value $\neg L(x)$, but the error probability is bounded above by $\epsilon(n)$. Then, for every $\epsilon:\mathbb N\to(0,1)$, we define a class $\mathsf{Heur}_\epsilon\mathsf P$ as a class of distributional problems that have a polynomial-time error-prone heuristic algorithm with failure probability $\epsilon(n)$. We also define an error-prone heuristic scheme and the class $\mathsf{HeurP}$ in the same manner as the errorless case.
We also define classes $\mathsf{AvgP/poly}$, $\mathsf{HeurP/poly}$, $\mathsf{AvgSIZE}[s(n)]$, and $\mathsf{HeurSIZE}[s(n)]$ for each size parameter $s(n)$ in the same manner as above.
\paragraph{Definition 11 (auxiliary-input pseudorandom generator).}
Let $n,m:\mathbb N\to\mathbb N$ be polynomials. We say that $G=\{G_z:\{0,1\}^{n(|z|)}\to\{0,1\}^{m(|z|)}\}_{z\in\{0,1\}^*}$ is an auxiliary-input pseudorandom generator (AIPRG) against $C$ if each $G_z(x)$ is polynomial-time computable from $(z,x)$, $n(\ell)<m(\ell)$ holds for any $\ell\in\mathbb N$, and for any adversary $A$ in $C$, there exists an infinite subset $Z_A\subseteq\{0,1\}^*$ such that for every $z\in Z_A$,
\[\left|\Pr[A(z,G_z(U_{n(|z|)}))=1]-\Pr[A(z,U_{m(|z|)})=1]\right|<\operatorname{negl}(|z|).\]
When the auxiliary-input $z$ is obvious in context, we write $n(|z|)$ and $m(|z|)$ as $n$ and $m$, respectively.
\paragraph{Definition 13 (hitting set generator).}
Let $\ell,m:\mathbb N\to\mathbb N$ be polynomials. We say that $G=\{G_n\}_{n\in\mathbb N}$, where $G_n:\{0,1\}^{\ell(n)}\to\{0,1\}^{m(n)}$ is a hitting set generator (HSG) against $C$ if each $G_n$ is polynomial-time computable, $\ell(n)<m(n)$ holds for each $n\in\mathbb N$, and $G$ hits any language recognized by adversaries in $C$ in the following sense: For any adversary $A$ in $C$, let $L_A=\{L_{A,n}\}_{n\in\mathbb N}$ be a language recognized by $A$, where $L_{A,n}\subseteq\{0,1\}^n$ for each $n\in\mathbb N$. Then, for infinitely many $n\in\mathbb N$, the following holds:
\[|L_{A,m(n)}|>2^{m(n)-1}\ \Longrightarrow\ L_{A,m(n)}\cap\operatorname{Im}G_n\neq\varnothing.\]

\paragraph{Theorem 4 (core oracle separation clauses).}
For any constant $a>0$, there exists an oracle $O$ relative to which the following hold:
\begin{itemize}
\item (Error-prone average-case easiness of $\mathsf{NP}$) $\mathsf{DistNP}^{O}\subseteq\mathsf{HeurP}^{O}$.
\item (Errorless average-case hardness of $\mathsf{NP}$) $\mathsf{DistNP}^{O}\not\subseteq\mathsf{AvgSIZE}^{O}[2^{an/\log n}]$.
\end{itemize}

\paragraph{4. Oracle Construction.}
In this section, we formally present the proposed oracle construction. Let $a>0$ be a parameter.
\paragraph{Construction.}
$\mathcal O_a=\mathcal F+\mathcal A$, where each oracle is randomly selected according to the following process:
\begin{enumerate}
\item Let $t(n)=2^{an/\log n}$ be the upper bound on the time of nonuniform adversaries and $c=7a$.
\item Define functions $p$ and $i_{\max}$ as $p(n)=t(n)^{-6}$ and $i_{\max}(n)=\frac1c\log\log t(n)$. Here, $p$ is a parameter of random restriction, and $i_{\max}$ is the number of applications of random restrictions.
\item For each $n\in\mathbb N$, define a set $V_{n,0}$ of variables taking binary values as follows:
\[V_{n,0}=\{F_{z,x,\ell}:z,x\in\{0,1\}^n,\ell\in[n]\}.\]
\item For each $n\in\mathbb N$, let $S_{n,0}=\{0,1\}^n$.
\item For each $n\in\mathbb N$ and $i\in[i_{\max}(n)-1]$, we inductively (on $i$) select a $p(n)$-random restriction $\rho^*_{n,i}$ to $V_{n,i-1}$ and a random subset $S_{n,i}\subseteq S_{n,i-1}$ of size $p(n)\cdot|S_{n,i-1}|$. Then, we define a restriction $\rho_{n,i}$ to $V_{n,i-1}$ and a subset $V_{n,i}\subseteq V_{n,i-1}$ as follows:
\[\rho_{n,i}(z,x,\ell)=\begin{cases}*&\text{if }z\in S_{n,i},\\\rho^*_{n,i}(z,x,\ell)&\text{otherwise},\end{cases}\]
\[V_{n,i}=\rho_{n,i}^{-1}(*)\left(=\rho_{n,i}^{*-1}(*)\cup\{F_{z,x,\ell}:z\in S_{n,i}\}\right).\]
We also define $\rho_{n,i_{\max}(n)}$ as a full assignment to $V_{n,i_{\max}(n)-1}$ selected uniformly at random. Let $\rho_{n,i}\equiv\rho_{n,i_{\max}(n)}$ for each $i\geq i_{\max}(n)+1$. We use the notation $\rho_{n,\leq i}$ to represent the composite restriction $\rho_{n,1}\cdots\rho_{n,i}$ to $V_{n,0}$ for each $n$ and $i$.
\item Define $\mathcal F=\{F_n\}_{n\in\mathbb N}$, where $F_n:\{0,1\}^n\times\{0,1\}^n\times[n]\to\{0,1\}$, as $F_n(z,x,\ell)=\rho_{n,\leq i_{\max}(n)}(z,x,\ell)$.
\item Define $\mathcal A$ as follows: On input $(M,x,1^{T^{2^c}})$, where $M$ is a nondeterministic oracle machine, $x\in\{0,1\}^*$, and $T\in\mathbb N$, the oracle $\mathcal A(M,x,1^{T^{2^c}})$ returns 0 or 1 according to the following procedure:
\begin{enumerate}
\item Let $i_T:=\frac1c\log\log T$.
\item Construct a $T$-DNF $\phi$ on variables in $V_{n,0}$ representing the execution of $M^{\mathcal F+\mathcal A}(x)$ in $T$ steps, where the top-most OR corresponds to the nondeterminism on a possible choice of $\mathcal F$ (say, $F'$) and an accepting path of $M^{F'+\mathcal A}(x)$, and each term performs verification whether $M$'s at most $T$ queries (say, $(z_1,x_1,\ell_1),\ldots,(z_q,x_q,\ell_q)$ for some $q\leq T$) are consistent with the actual choices of $\mathcal F$, i.e., for each $i\in[q]$, the term contains $F_{z_i,x_i,\ell_i}$ if $F'(z_i,x_i,\ell_i)=1$; otherwise, $\neg F_{z_i,x_i,\ell_i}$ as a literal.
\item If $\phi|_{\rho_{1,\leq i_T},\ldots,\rho_{T,\leq i_T}}\equiv b$ for some $b\in\{0,1\}$, then return $b$, otherwise, return 0.
\end{enumerate}
\end{enumerate}
We can verify that $\mathcal A$ is well-defined (i.e., not circular on recursive calls for $\mathcal A$) as follows.
\paragraph{Proposition 14.}
For each input, the value of $\mathcal A(M,x,1^{T^{2^c}})$ is determined only by $\rho_{n,j}$ (equivalently, $\rho^*_{n,j}$ and $S_{n,j}$) for $n\leq T$ and $j\leq i_T$ (remember that $i_T=\frac1c\log\log T$).
\begin{proof}
We show the proposition by induction on $T$. Remember that, on input $(M,x,1^{T^{2^c}})$, the oracle $\mathcal A$ first makes a $T$-DNF $\phi$ based on $M$ independently of the values of $\mathcal F$.
Suppose that $M$ makes some query $(M',x',1^{T'^{2^c}})$ to $\mathcal A$ for constructing $\phi$. Since the length of such a query is at most $T$, we have $T'^{2^c}\leq T$ and
\[i_{T'}=\frac1c\log\log T'\leq\frac1c\log\log T^{1/2^c}=\frac1c\log\log T-1=i_T-1.\]
By the induction hypothesis, the answer of $\mathcal A(M',x',1^{T'^{2^c}})$ is determined by only $\rho_{n,j}$ for $n\leq T'$ and $j\leq i_T-1$, and so is $\phi$. Then, $\mathcal A$ determines the answer by restricting $\phi$ by $\rho_{n,j}$ for $n\leq T$ and $j\leq i_T$. Therefore, $\mathcal A(M,x,1^{T^{2^c}})$ is determined only by $\rho_{n,j}$ for $n\leq T$ and $j\leq i_T$.
\end{proof}
When the parameter $a$ is clear from the context, we omit the subscript $a$ from $\mathcal O_a$.

\paragraph{5. Error-Prone Average-Case Easiness of NP.}
In this section, we show the error-prone average-case easiness of $\mathsf{NP}$.
\paragraph{Theorem 15.}
With probability 1 over the choice of $\mathcal O_a$, $\mathsf{DistNP}^{\mathcal O_a}\subseteq\mathsf{HeurP}^{\mathcal O_a}$ holds.
First, we introduce several notations. For each choice of $\mathcal O$, we define the oracle $\mathcal A^*$ in the same manner as $\mathcal A$ except we apply $\rho_{1,\leq i_T-1}\rho^*_{1,i_T},\ldots,\rho_{T,\leq i_T-1}\rho^*_{T,i_T}$ to $\phi$ instead of $\rho_{1,\leq i_T},\ldots,\rho_{T,\leq i_T}$. Note that $\mathcal A^*$ executes a given nondeterministic machine $M$ with access to $\mathcal A$ (rather than $\mathcal A^*$) to construct the corresponding DNF $\phi$. We can verify that $\mathcal A^*$ is well-defined (i.e., not circular) in the same manner as Proposition 14.
Now, we show that $\mathcal A$ and $\mathcal A^*$ do not differ considerably.
\paragraph{Lemma 16.}
For each input $(M,x,1^{T^{2^c}})$ to $\mathcal A$, we have that
\[\Pr_{\mathcal O}[\mathcal A(M,x,1^{T^{2^c}})\neq\mathcal A^*(M,x,1^{T^{2^c}})]=O(T^{-4}).\]
\begin{proof}
Let $i:=i_T=(1/c)\log\log T$. For all $n\leq T$ and $j\leq i-1$, we fix $\rho^*_{n,j}$, $S_{n,j}$, and $\rho^*_{n,i}$ arbitrarily; let $C_T$ denote this condition. Notice that the DNF formula $\phi_{C_T}$ constructed by $\mathcal A$ and $\mathcal A^*$ is determined only by $C_T$ because all answers to recursive calls for $\mathcal A$ are determined by $C_T$ as in Proposition 14. Let $\phi'_{C_T}=\phi_{C_T}|_{\rho_{1,\leq i-1},\ldots,\rho_{T,\leq i-1}}$. Then, the value of $\mathcal A(M,x,1^{T^{2^c}})$ (resp. $\mathcal A^*(M,x,1^{T^{2^c}})$) is determined by $\phi'_{C_T}|_{\rho_{1,i},\ldots,\rho_{T,i}}$ (resp. $\phi'_{C_T}|_{\rho^*_{1,i},\ldots,\rho^*_{T,i}}$).
For any DNF formula $\phi$ and a restriction $\rho$, there are the following three cases: (i) $\phi|_\rho\equiv0$, (ii) $\phi|_\rho\equiv1$, or (iii) $\phi|_\rho$ does not become a constant (we write this case as $\phi|_\rho\equiv*$). Following this case analysis, there exist $3^2=9$ cases on $(\phi'_{C_T}|_{\rho^*_{1,i},\ldots,\rho^*_{T,i}},\phi'_{C_T}|_{\rho_{1,i},\ldots,\rho_{T,i}})$. However, since each $\rho_{n,i}$ is a subrestriction of $\rho^*_{n,i}$ (i.e., $\rho_{n,i}$ assigns values only to variables that are also assigned by $\rho^*_{n,i}$), the following 4 cases do not occur: $(0,1)$, $(1,0)$, $(*,0)$, and $(*,1)$. Further, in the cases of $(0,0)$, $(1,1)$, $(*,*)$, and $(*,0)$, the answers by $\mathcal A^*$ and $\mathcal A$ do not differ because $\mathcal A^*$ and $\mathcal A$ return 0 in the case of $*$. Thus, we only need to consider the case of $(1,*)$.
By the aforementioned argument, the probability in the lemma is expressed as follows:
\begin{align*}
&\Pr_{\mathcal O}[\mathcal A(M,x,1^{T^{2^c}})\neq\mathcal A^*(M,x,1^{T^{2^c}})]\\
&=\mathbb E_{C_T}\!\left[\Pr_{S_{1,i},\ldots,S_{T,i}}\!\left[\phi'_{C_T}|_{\rho_{1,i},\ldots,\rho_{T,i}}\equiv*\ \middle|\ C_T,\phi'_{C_T}|_{\rho^*_{1,i},\ldots,\rho^*_{T,i}}\equiv1\right]\right].
\end{align*}
To bound the probability in the right-hand side for each condition, we consider the case in which $\phi'_{C_T}|_{\rho^*_{1,i},\ldots,\rho^*_{T,i}}\equiv1$. Then, we can select a term $\tau$ in $\phi'_{C_T}$ such that $\tau|_{\rho^*_{1,i},\ldots,\rho^*_{T,i}}\equiv1$. For each $n\leq T$, define $Z_n\subseteq\{0,1\}^n$ as
\[Z_n=\{z\in\{0,1\}^n:\exists(x,\ell)\in\{0,1\}^n\times[n]\text{ s.t. a variable }F_{z,x,\ell}\text{ is contained in }\tau\}.\]
Since $\phi'_{C_T}$ is a $T$-DNF, $|Z_n|\leq T$ for each $n\leq T$. Furthermore, for any $n\in\mathbb N$ such that $i_{\max}(n)\leq i-1$, we have that $Z_n=\varnothing$ because all such variables must be assigned by $\rho_{n,\leq i-1}$. If $\phi'_{C_T}|_{\rho_{1,i},\ldots,\rho_{T,i}}\equiv*$, then $\tau|_{\rho_{1,i},\ldots,\rho_{T,i}}\equiv*$ must hold. This event occurs only if $\bigcup_{n\leq T}(S_{n,i}\cap Z_n)\neq\varnothing$ holds. Since each $S_{n,i}$ is selected from $S_{n,i-1}$ uniformly at random, this probability is bounded above by
\begin{align*}
\Pr\left[\bigcup_{n\leq T}(S_{n,i}\cap Z_n)\neq\varnothing\right]
&\leq\sum_{n\leq T}\Pr[S_{n,i}\cap Z_n\neq\varnothing]\\
&\leq\sum_{n\leq T:i_{\max}(n)\geq i}|Z_n|\cdot|S_{n,i}|/|S_{n,i-1}|\\
&\leq O(T)\cdot\sum_{n:t^{-1}(T)\leq n\leq T}p(n)\\
&\leq O(T^2\cdot t(t^{-1}(T))^{-6})=O(T^{-4}).
\end{align*}
Thus, we conclude that
\begin{align*}
&\Pr_{\mathcal O}[\mathcal A(M,x,1^{T^{2^c}})\neq\mathcal A^*(M,x,1^{T^{2^c}})]\\
&=\mathbb E_{C_T}\!\left[\Pr_{S_{1,i},\ldots,S_{T,i}}\!\left[\phi'_{C_T}|_{\rho_{1,i},\ldots,\rho_{T,i}}\equiv*\ \middle|\ C_T,\phi'_{C_T}|_{\rho^*_{1,i},\ldots,\rho^*_{T,i}}\equiv1\right]\right]\leq O(T^{-4}).
\end{align*}
\end{proof}
Furthermore, we can show the average-case easiness of $\mathsf{NP}$ under the oracle access to $\mathcal A^*$ (instead of $\mathcal A$). This part is similar to in [25,20]. Thus, we defer the proof to Appendix A.
\paragraph{Lemma 17 ([25,20]).}
Let $M$ be a $t_M(n)$-time nondeterministic oracle machine and $S$ be a randomized polynomial-time oracle sampling machine. We assume that $S(1^n)$ takes at most $t_S(n)$ time to generate an instance of length $n$. Then, the following event occurs with probability 1 over the choice of $\mathcal O$: for any $n\in\mathbb N$,
\[\Pr_{x\leftarrow S^{\mathcal O}(1^n)}[M^{\mathcal O}(x)\neq\mathcal A^*(M,x,1^{T^{2^c}})]\leq O(n^{-4}),\]
where $T=\max\{n^{2^c},t_M(n)^{2^c},t_S(n)^{2^c}\}$.
Now, we derive the average-case easiness of $\mathsf{NP}$ from Lemmas 16 and 17.
\begin{proof}[Proof of Theorem 15]
We consider an arbitrary distributional $\mathsf{NP}$ problem $(L,D)$ and assume that $L$ is specified by a $t_M(n)$-time nondeterministic oracle machine $M$, and $D$ is specified by a randomized $t_S(n)$-time oracle sampling machine $S$. Let $T=\max\{n^{2^c},t_M(n)^{2^c},t_S(n)^{2^c}\}$ as in Lemma 17.
We construct an error-prone heuristic algorithm $B$ for $(L,D)$ as follows: On input $x\in\{0,1\}^n$, $B$ queries $(M,x,1^{T^{2^c}})$ to $\mathcal A$ and returns the same answer. In the following, we will verify that the error probability of $B$ (over the choice of $\mathcal O$ and $S^{\mathcal O}(1^n)$) is bounded above by $O(n^{-4})$ for each input size $n$. Then, by applying Markov's inequality and the Borel--Cantelli lemma, the error probability of $B$ is bounded above by $n^{-2}$ for all sufficiently large $n$ with probability 1 over the choice of $\mathcal O$. Since the number of tuples $(M,S)$ is countable, we can conclude that all distributional $\mathsf{NP}$ problems have an error-prone heuristic algorithm with error probability at most $n^{-2}$ with probability 1 over the choice of $\mathcal O$. Based on the argument in [24, Proposition 3], this is sufficient for the statement that $\mathsf{DistNP}^{\mathcal O}\subseteq\mathsf{HeurP}^{\mathcal O}$.
Therefore, it is sufficient to show that, for any $n\in\mathbb N$,
\[\Pr_{\mathcal O,x\leftarrow S^{\mathcal O}(1^n)}[M^{\mathcal O}(x)\neq\mathcal A(M,x,1^{T^{2^c}})]\leq O(n^{-4}).\]
Obviously, this event occurs only if (i) $\mathcal A(M,x,1^{T^{2^c}})\neq\mathcal A^*(M,x,1^{T^{2^c}})$ or (ii) $M^{\mathcal O}(x)\neq\mathcal A^*(M,x,1^{T^{2^c}})$ occur. By Lemmas 16 and 17 and the union bound, we have
\begin{align*}
&\Pr_{\mathcal O,x\leftarrow S^{\mathcal O}(1^n)}[M^{\mathcal O}(x)\neq\mathcal A(M,x,1^{T^{2^c}})]\\
&\leq\Pr_{\mathcal O,x\leftarrow S^{\mathcal O}(1^n)}[\mathcal A(M,x,1^{T^{2^c}})\neq\mathcal A^*(M,x,1^{T^{2^c}})]\\
&\quad+\Pr_{\mathcal O,x\leftarrow S^{\mathcal O}(1^n)}[M^{\mathcal O}(x)\neq\mathcal A^*(M,x,1^{T^{2^c}})]\\
&=O(T^{-4})+O(n^{-4})=O(n^{-4})+O(n^{-4})=O(n^{-4}).
\end{align*}
\end{proof}

\paragraph{A. Proof of Lemma 17.}
Fix $n\in\mathbb N$ arbitrarily. Remember that $T=\max\{n^{2^c},t_M(n)^{2^c},t_S(n)^{2^c}\}$. Let $i_T=c^{-1}\log\log T$.
We first show that the instance $x\in\{0,1\}^n$ generated by $S^{\mathcal O}(1^n)$ is determined by only $\rho_{n',j}$ for $n'\leq T$ and $j\leq i_T-1$ with probability at least $1-O(n^{-5})$. Then, we show that $\mathcal A^*(M,x,1^{T^{2^c}})$ returns $M^{\mathcal O}(x)$ with probability at least $1-O(n^{-5})$ under the condition that $x$ is determined by only $\rho_{n',j}$ for $n'\leq T$ and $j\leq i_T-1$. If we assume these, then by the union bound, we have the lemma as
\[\Pr_{\mathcal O,x\leftarrow S^{\mathcal O}(1^n)}[M^{\mathcal O}(x)\neq\mathcal A^*(M,x,1^{T^{2^c}})]\leq O(n^{-5})+O(n^{-5})=O(n^{-5})\ (=O(n^{-4})).\]
Now, we show the first claim. Since $t_S(n)\leq T^{1/2^c}$, the answers of $\mathcal A$ to queries made by $S^{\mathcal O}(1^n)$ are determined by only $\rho_{n',j}$ for $n'\leq T^{1/2^c}$ and $j\leq i_T-2$. Under an arbitrary condition on restrictions $\rho_{n',j}$ for $n'\leq T^{1/2^c}$ and $j\leq i_{t_S(n)^{1/2^c}}\leq i_T-2$, the output $S^{\mathcal O}(1^n)$ is determined by only $\rho_{n',j}$ for $n'\leq T$ and $j\leq i_T-1$ unless $S$ queries $(z,x,\ell)\in\bigcup_{n'\leq T^{1/2^c}}\rho_{n',i_T-1}^{*-1}(*)$ to $\mathcal F$. Note that, if $n'\in\mathbb N$ satisfies $n'<t^{-1}(T^{1/2^c})$, then we have $i_{\max}(n')<c^{-1}\log\log(T^{1/2^c})=c^{-1}\log\log T-1=i_T-1$. Thus, $\rho_{n',i_T-1}^{*-1}(*)=\varnothing$. Otherwise, each element in $\rho_{n',i_T-1}^{*-1}(*)$ is selected from $V_{n',i_T-2}$ independently with probability $p(n')$. Since $S^{\mathcal O}(1^n)$ accesses to $\mathcal A$ at most $T^{1/2^c}$ times, such a conditional probability is bounded above by
\begin{align*}
T^{1/2^c}\max_{t^{-1}(T^{1/2^c})\leq n'\leq T^{1/2^c}}p(n')
&=T^{1/2^c}p(t^{-1}(T^{1/2^c}))\\
&=T^{1/2^c}t(t^{-1}(T^{1/2^c}))^{-6}\\
&=(T^{1/2^c})^{-5}\leq n^{-5},
\end{align*}
where the inequality holds because $T\geq n^{2^c}$.
Next, we show the second claim. Under the condition that the given instance $x\in\{0,1\}^n$ is determined by only $\rho_{n',j}$ for $n'\leq T$ and $j\leq i_T-1$, the $T$-DNF formula $\phi$ constructed in $\mathcal A^*(M,x,1^{T^{2^c}})$ is determined only by $\rho_{n',j}$ for $n'\leq T$ and $j\leq i_T-1$. Then, applying the restriction $\rho_{n',j}$ for $n'\leq T$ and $j\leq i_T$ under this condition is regarded as a $p(n')$-random restriction to $V_{n',i_T-1}$ for each $n'\leq T$, where we can ignore small $n'$ such that $n'<t^{-1}(T)$ because $i_{\max}(n')<c^{-1}\log\log T=i_T$ for such $n'$. By applying the switching lemma (particularly, the baby switching lemma) for $T$-DNF [14], the probability that $\phi$ does not become a constant is at most
\begin{align*}
O\left(T\max_{t^{-1}(T)\leq n'\leq T}p(n')\right)&=O(T\cdot t(t^{-1}(T))^{-6})\\
&=O(T^{-5})=O(n^{-5}),
\end{align*}
where the last equation holds because $T\geq n^{2^c}\geq n$. Remember that $\mathcal A^*$ always returns the correct answer whenever $\phi$ becomes constant. Therefore, the second claim holds.

\paragraph{6. Errorless Average-Case Hardness of NP.}
In this section, we show the hardness part of our oracle separation. First, we show the existence of AIPRG relative to $\mathcal O_a$. Then we show the other hardness results, including the errorless average-case hardness for $\mathsf{NP}$, as corollaries.
We use the following theorem, which shows the existence of PRGs based on a random oracle.
\paragraph{Theorem 18 ([45]).}
For each $n\in\mathbb N$, let $R_n:\{0,1\}^n\to\{0,1\}^n$ be a random function oracle, i.e., $R_n$ is selected uniformly at random from $\{f:\{0,1\}^n\to\{0,1\}^n\}$.
There exist a polynomial-time deterministic oracle machine $G^?$ and constants $c\geq1$ and $b,\epsilon>0$ satisfying the following: For any $n\in\mathbb N$ and $x\in\{0,1\}^{cn}$, $G^{R_n}(x)$ generates a binary string of length $4cn$, and all oracle circuits $C^?$ of size $2^{bn}$ satisfy that
\[\left|\Pr_{U_{cn}}[C^{R_n}(G^{R_n}(U_{cn}))=1]-\Pr_{U_{4cn}}[C^{R_n}(U_{4cn})=1]\right|\leq2^{-\epsilon n},\]
with probability at least $1-2^n$ over the choice of $R_n$.
Furthermore, the result above is relativized, i.e., the above holds in the presence of an arbitrary oracle $O$ independent of the choice of $R_n$.
Now, we show the existence of AIPRG relative to $\mathcal O_a$.
\paragraph{Theorem 19.}
Let $a>0$ be an arbitrary constant. With probability 1 over the choice of $\mathcal O_a$, there exists an AIPRG $G^{\mathcal O_a}=\{G_z^{\mathcal O_a}\}_{z\in\{0,1\}^*}$ against $\mathsf{SIZE}^{\mathcal O_a}[2^{\epsilon an/\log n}]$ for some absolute constant $\epsilon>0$, where $G_z^{\mathcal O_a}:\{0,1\}^{|z|}\to\{0,1\}^{3|z|}$ for each $z\in\{0,1\}^*$.
\begin{proof}
Let $G^?$ and $c$ be the oracle machine and the constant in Theorem 18, respectively. Then, we define an AIPRG $G'=\{G'_z\}_{z\in\{0,1\}^*}$ as $G_z'^{\mathcal O}(x)=G^{F_{z'}}(x')_{\leq3|z|}$, where $x\in\{0,1\}^{|z|}$, $z'=z_{\leq\lfloor|z|/c\rfloor}$, $x'=x_{\leq c|z'|}$, and $F_{z'}:\{0,1\}^{|z'|}\to\{0,1\}^{|z'|}$ is defined as $F_{z'}(y)=\mathcal F(z',y,1)\circ\cdots\circ\mathcal F(z',y,|z'|)$. The validity of the truncation is verified as that $|x'|=c|z'|\leq c\cdot|z|/c=|z|=|x|$ and $|G^{F_{z'}}(x')|=4|x'|=4c|z'|\geq4c(|z|/c-1)\geq4|z|-4c\geq3|z|$ for any $z$ with $|z|\geq4c$.
Let $\epsilon=1/2c$ and $s(n)=2^{\epsilon an/\log n}$. We show that $G'$ above is an AIPRG against $\mathsf{SIZE}^{\mathcal O}[s(n)]$. Suppose there exists a family $C=\{C_n\}_{n\in\mathbb N}$ of oracle circuits of size $s(n)$ that breaks $G'$, i.e., for any sufficiently large $n\in\mathbb N$ and any $z\in\{0,1\}^n$,
\[\left|\Pr_{U_n}[C_n^{\mathcal O}(z,G_z'^{\mathcal O}(U_n))=1]-\Pr_{U_{3n}}[C_n^{\mathcal O}(z,U_{3n})=1]\right|>\frac1{\operatorname{poly}(n)}.\]
Fix $n\in\mathbb N$ arbitrarily, and let $n'=\lfloor n/c\rfloor$. We consider an arbitrary choice of $\mathcal O$ except for the values of $\rho_{n',i_{\max}(n')}$ (we write this condition as $R$ for convenience). Fix $z\in\{0,1\}^n$ such that $z'=z_{\leq n'}\in S_{n',i_{\max}(n')-1}$ arbitrary (where $S_{\cdot,\cdot}$ is the random set in the oracle construction). We refer to such $z$ as a hard index.
Since the size of $C_n$ is at most $s(n)=2^{an/(2c\log n)}\leq2^{a\lfloor n/c\rfloor/\log\lfloor n/c\rfloor}=t(n')$ for sufficiently large $n$ (where $t$ is the time-bound function in the oracle construction), the answers to queries by $C_n$ of the form $\mathcal A(M,y,1^{T^{2^c}})$ do not depend on the values of $F_{z'}$ and are determined by condition $R$ because they are determined only by the restrictions $\rho_{\cdot,j}$ for
\[j\leq c^{-1}\log\log T\leq c^{-1}\log\log s(n)^{1/2^c}\leq c\log\log t(n')-1<i_{\max}(n').\]
Now, we consider an arbitrary choice of $\rho_{n',i_{\max}(n')}$ except for the values of $F_{z'}$ and denote this condition by $R'$. It is not hard to verify that $F_{z'}$ is selected uniformly at random even under the conditions $R$ and $R'$. Therefore, under the conditions $R$ and $R'$, we can identify the query access to $\mathcal O$ by $C$ with the query access to another oracle $O'$ (determined only by $R$ and $R'$) and a random function oracle $F_{z'}$ that are selected independently of $O'$.
For any $n\in\mathbb N$, let $E_n$ be an event (over the choice of $F_{z'}$) that there exists a circuit $C'$ of size $2^{bn'}$, where $b$ represents the constant in Theorem 18, such that
\begin{align*}
&\left|\Pr_{U_n}[C'^{\mathcal O}(G_z'^{F_{z'}}(U_n))=1]-\Pr_{U_{3n}}[C'^{\mathcal O}(U_{3n})=1]\right|\\
&=\left|\Pr_{U_n}[C'^{O',F_{z'}}(G^{F_{z'}}(U_n))=1]-\Pr_{U_{3n}}[C'^{O',F_{z'}}(U_{3n})=1]\right|>\frac1{\operatorname{poly}(n)}.
\end{align*}
Then, by Theorem 18 (relative to $O'$), we have $\Pr_{\mathcal O}[E_n\mid R,R']\leq2^{-\Omega(n)}$. By the Borel--Cantelli lemma, $E_n$ occurs only for finitely many $n$'s with probability 1 over the choice of $\mathcal O$ conditioned on $R,R'$ (i.e., the choice of $F_{z'}$). By taking the expectation over $R,R'$, we can show that, with probability 1 over the choice of $\mathcal O$, there is no family $C'$ of $2^{bn'}$-size circuits satisfying that for any sufficiently large $n\in\mathbb N$, there exists a hard index $z\in\{0,1\}^n$ such that
\[\left|\Pr_{U_n}[C'^{\mathcal O}(G_z'^{F_{z'}}(U_n))=1]-\Pr_{U_{3n}}[C'^{\mathcal O}(U_{3n})=1]\right|>\frac1{\operatorname{poly}(n)}.\]
However, the circuit $C$ in the assumption violates this statement by embedding a hard index $z$ as auxiliary-input because the size is at most $n+s(n)=O(2^{(a/2c)\cdot n/\log n})=o(2^{bn'})$. Therefore, we conclude that, with probability 1 over the choice of $\mathcal O$, there is no such a circuit $C$ of size $s(n)$, and $G'$ is an AIPRG against $\mathsf{SIZE}[s(n)]$.
\end{proof}
Next, AIPRG implies HSG by regarding the auxiliary-input as a part of the hidden input to HSG. It is not hard to verify the security (refer to [34, Lemma 18] for the formal proof).
\paragraph{Corollary 21.}
Let $a>0$ be an arbitrary constant. With probability 1 over the choice of $\mathcal O_a$, there exists an HSG $G^{\mathcal O_a}=\{G_n^{\mathcal O_a}\}_{n\in\mathbb N}$ against $\mathsf{SIZE}^{\mathcal O_a}[2^{\epsilon an/\log n}]$ for some absolute constant $\epsilon>0$, where $G_n^{\mathcal O_a}:\{0,1\}^{2n}\to\{0,1\}^{3n}$ for each $n\in\mathbb N$.
Furthermore, the existence of HSGs implies the errorless average-case hardness of $\mathsf{NP}$, as observed in [21] in the context of natural proofs and the average-case complexity of the minimum circuit size problem.
\paragraph{Corollary 22.}
Let $a>0$ be an arbitrary constant. With probability 1 over the choice of $\mathcal O_a$, $\mathsf{DistNP}^{\mathcal O_a}\not\subseteq\mathsf{AvgSIZE}^{\mathcal O_a}[2^{\epsilon an/\log n}]$ for some absolute constant $\epsilon>0$.
\begin{proof}
Let $G^{\mathcal O}$ and $\epsilon$ be the HSG and the constant in Corollary 21, respectively. Then, we define the language $L^{\mathcal O}$ as $L^{\mathcal O}:=\operatorname{Im}G^{\mathcal O}$. Obviously, $L^{\mathcal O}\in\mathsf{NP}^{\mathcal O}$ and $(L^{\mathcal O},\{U_n\}_{n\in\mathbb N})\in\mathsf{DistNP}^{\mathcal O}$. Thus, it is sufficient to show that $(L^{\mathcal O},\{U_n\}_{n\in\mathbb N})\notin\mathsf{Avg}_{1/4}\mathsf{SIZE}^{\mathcal O}[2^{\epsilon an/(4\log n)}]$.
For contradiction, we assume that $(L^{\mathcal O},\{U_n\}_{n\in\mathbb N})\in\mathsf{Avg}_{1/4}\mathsf{SIZE}^{\mathcal O}[2^{\epsilon an/(4\log n)}]$. Then, there exists a family $C=\{C_n\}_{n\in\mathbb N}$ of $O(2^{\epsilon an/(4\log n)})$-size oracle circuits for $(L^{\mathcal O},\{U_n\}_{n\in\mathbb N})$, i.e., for any sufficiently large $n\in\mathbb N$,
\[\Pr_{y\leftarrow_u\{0,1\}^{3n}}[C_{3n}^{\mathcal O}(y)=\bot]\leq1/4\]
and for each $y\in\{0,1\}^{3n}$,
\[C_{3n}^{\mathcal O}(y)\in\{\mathbf1(y\in L^{\mathcal O}),\bot\}.\]
Note that the size of $C_{3n}$ is at most $O(2^{\epsilon a\cdot3n/(4\log3n)})$.
Next, we define an adversary $C'$ for $G^{\mathcal O}$ as follows: for a given $y\in\{0,1\}^{3n}$ (i.e., the length of seed is $2n$), $C_{2n}'^{\mathcal O}$ simulates $C_{3n}^{\mathcal O}(y)$, and if $C_{3n}$ returns 1 or $\bot$, then $C'_{2n}$ outputs 0; otherwise (i.e., if $C_{3n}$ returns 0), $C'_{2n}$ outputs 1. Then, based on the aforementioned inequalities, it is not hard to verify that for any sufficiently large $n\in\mathbb N$,
\begin{align*}
\Pr_{y\leftarrow_u\{0,1\}^{3n}}[C_{2n}'^{\mathcal O}(y)=0]
&\leq\frac14+\frac{|\{G^{\mathcal O}(x):x\in\{0,1\}^{2n}\}|}{|\{0,1\}^{3n}|}\\
&\leq\frac14+\frac{2^{2n}}{2^{3n}}<\frac12,
\end{align*}
and for any $x\in\{0,1\}^{2n}$, we have $C_{3n}^{\mathcal O}(G^{\mathcal O}(x))\in\{1,\bot\}$ and $C_{2n}'^{\mathcal O}(G^{\mathcal O}(x))=0$.
Therefore, $C'$ succeeds in avoiding $\operatorname{Im}G^{\mathcal O}$, and the size is bounded above by $O(2^{\epsilon an/\log n})=O(2^{\epsilon a(2n)/\log(2n)})$. This contradicts Corollary 21. Thus, we conclude that $(L^{\mathcal O},\{U_n\}_{n\in\mathbb N})\notin\mathsf{Avg}_{1/4}\mathsf{SIZE}^{\mathcal O}[2^{\epsilon an/(4\log n)}]$.
\end{proof}

\begin{thebibliography}{99}
\bibitem[1]{published-1} S Aaronson and A Wigderson. Algebrization: A New Barrier in Complexity Theory. ACM Trans. Comput. Theory, 1(1), February 2009.
\bibitem[2]{published-2} E. Allender, M. Cheraghchi, D. Myrisiotis, H. Tirumala, and I. Volkovich. One-Way Functions and a Conditional Variant of MKTP. In 41st IARCS Annual Conference on Foundations of Software Technology and Theoretical Computer Science, FSTTCS 2021, December 15-17, 2021, Virtual Conference, volume 213 of LIPIcs, pages 7:1–7:19. Schloss Dagstuhl - Leibniz-Zentrum für Informatik, 2021.
\bibitem[3]{published-3} B. Applebaum, B. Barak, and D. Xiao. On Basing Lower-Bounds for Learning on Worst-Case Assumptions. In Proceedings of the 49th Annual IEEE Symposium on Foundations of Computer Science, FOCS’08, pages 211–220, 2008.
\bibitem[4]{published-4} T. Baker, J. Gill, and R. Solovay. Relativizations of the P =?N P Question. SIAM Journal on Computing, 4(4):431–442, 1975.
\bibitem[5]{published-5} Manuel Blum and Sampath Kannan. Designing Programs that Check Their Work. J. ACM, 42(1):269–291, 1995. doi:10.1145/200836.200880.
\bibitem[6]{published-6} Manuel Blum and Silvio Micali. How to Generate Cryptographically Strong Sequences of Pseudo-Random Bits. SIAM J. Comput., 13(4):850–864, 1984. doi:10.1137/0213053.
\bibitem[7]{published-7} A. Bogdanov and L. Trevisan. Average-Case Complexity. Foundations and Trends in Theoretical Computer Science, 2(1):1–106, 2006.
\bibitem[8]{published-8} G. Brassard. Relativized Cryptography. IEEE Transactions on Information Theory, 29(6):877– 894, 1983.
\bibitem[9]{published-9} Harry Buhrman, Lance Fortnow, and Aduri Pavan. Some Results on Derandomization. Theory Comput. Syst., 38(2):211–227, 2005. doi:10.1007/s00224-004-1194-y.
\bibitem[10]{published-10} M. Carmosino, R. Impagliazzo, V. Kabanets, and A. Kolokolova. Learning Algorithms from Natural Proofs. In Proceedings of the 31st Conference on Computational Complexity, CCC’16. Schloss Dagstuhl–Leibniz-Zentrum fuer Informatik, 2016.
\bibitem[11]{published-11} R. Gennaro and L. Trevisan. Lower bounds on the efficiency of generic cryptographic constructions. In Proceedings 41st Annual Symposium on Foundations of Computer Science, pages 305–313, 2000.
\bibitem[12]{published-12} O. Goldreich. Foundations of Cryptography: Volume 1. Cambridge University Press, New York, NY, USA, 2006.
\bibitem[13]{published-13} O. Goldreich, S. Goldwasser, and S. Micali. How to Construct Random Functions. J. ACM, 33(4):792–807, August 1986.
\bibitem[14]{published-14} J. Håstad. Computational Limitations of Small-depth Circuits. MIT Press, Cambridge, MA, USA, 1987.
\bibitem[15]{published-15} J. Håstad, R. Impagliazzo, L. A. Levin, and M. Luby. A Pseudorandom Generator from Any One-way Function. SIAM J. Comput., 28(4):1364–1396, March 1999.
\bibitem[16]{published-16} S. Hirahara. Non-Black-Box Worst-Case to Average-Case Reductions within NP. In 59th IEEE Annual Symposium on Foundations of Computer Science, FOCS 2018, Paris, France, October 7-9, 2018, pages 247–258, 2018.
\bibitem[17]{published-17} S Hirahara. Characterizing Average-Case Complexity of PH by Worst-Case Meta-Complexity. In IEEE 61st Annual Symposium on Foundations of Computer Science (FOCS 2020), pages 50–60, 2020.
\bibitem[18]{published-18} S. Hirahara. Non-Disjoint Promise Problems from Meta-Computational View of Pseudorandom Generator Constructions. In 35th Computational Complexity Conference (CCC 2020), volume 169 of LIPIcs, pages 20:1–20:47, Dagstuhl, Germany, 2020.
\bibitem[19]{published-19} S. Hirahara. Average-Case Hardness of NP from Exponential Worst-Case Hardness Assump- tions. In 53rd Annual ACM Symposium on Theory of Computing (STOC 2021), 2021.
\bibitem[20]{published-20} S. Hirahara and M. Nanashima. On Worst-Case Learning in Relativized Heuristica. In 62nd Annual IEEE Symposium on Foundations of Computer Science (FOCS 2021), 2021.
\bibitem[21]{published-21} S. Hirahara and R. Santhanam. On the Average-Case Complexity of MCSP and Its Variants. In Ryan O’Donnell, editor, 32nd Computational Complexity Conference, CCC 2017, July 6-9, 2017, Riga, Latvia, volume 79 of LIPIcs, pages 7:1–7:20. Schloss Dagstuhl - Leibniz-Zentrum für Informatik, 2017.
\bibitem[22]{published-22} S. Hirahara and R. Santhanam. Errorless versus Error-prone Average-case Complexity. In The 13th Innovations in Theoretical Computer Science (ITCS 2022), 2022.
\bibitem[23]{published-23} R. Ilango, H. Ren, and R. Santhanam. Hardness on any Samplable Distribution Suffices: New Characterizations of One-Way Functions by Meta-Complexity. Electron. Colloquium Comput. Complex., page 82, 2021.
\bibitem[24]{published-24} R. Impagliazzo. A personal view of average-case complexity. In Proceedings of IEEE Tenth Annual Conference on Structure in Complexity Theory, pages 134–147, 1995.
\bibitem[25]{published-25} R. Impagliazzo. Relativized Separations of Worst-Case and Average-Case Complexities for NP. In 2011 IEEE 26th Annual Conference on Computational Complexity, pages 104–114, 2011.
\bibitem[26]{published-26} R. Impagliazzo and L. Levin. No better ways to generate hard NP instances than picking uniformly at random. In Proceedings of the 31st Annual Symposium on Foundations of Computer Science, FOCS’90, pages 812–821, 1990.
\bibitem[27]{published-27} R. Impagliazzo and S. Rudich. Limits on the Provable Consequences of One-Way Permutations. In Proceedings of the Twenty-First Annual ACM Symposium on Theory of Computing, STOC ’89, pages 44–61, New York, NY, USA, 1989. ACM.
\bibitem[28]{published-28} K. Ko. On the Complexity of Learning Minimum Time-Bounded Turing Machines. SIAM Journal on Computing, 20(5):962–986, 1991.
\bibitem[29]{published-29} L. A Levin. Average Case Complete Problems. SIAM J. Comput., 15(1):285–286, February 1986.
\bibitem[30]{published-30} Y. Liu and R. Pass. On One-way Functions and Kolmogorov Complexity. In IEEE 61st Annual Symposium on Foundations of Computer Science (FOCS 2020), pages 1243–1254, 2020. doi:10.1109/FOCS46700.2020.00118.
\bibitem[31]{published-31} Y. Liu and R. Pass. A Note on One-way Functions and Sparse Languages. Electron. Colloquium Comput. Complex., page 92, 2021.
\bibitem[32]{published-32} Y. Liu and R. Pass. On the Possibility of Basing Cryptography on EXP $\neq$ BPP. In Advances in Cryptology - CRYPTO 2021 - 41st Annual International Cryptology Conference, CRYPTO 2021, Virtual Event, August 16-20, 2021, Proceedings, Part I, volume 12825 of Lecture Notes in Computer Science, pages 11–40. Springer, 2021.
\bibitem[33]{published-33} M. Nanashima. A Theory of Heuristic Learnability. In Proceedings of the 34th Conference on Learning Theory, COLT’21. PMLR, August 2021.
\bibitem[34]{published-34} M. Nanashima. On Basing Auxiliary-Input Cryptography on NP-Hardness via Nonadaptive Black-Box Reductions. In James R. Lee, editor, 12th Innovations in Theoretical Computer Science Conference (ITCS 2021), volume 185 of LIPIcs, pages 29:1–29:15, Dagstuhl, Germany, 2021. Schloss Dagstuhl–Leibniz-Zentrum für Informatik.
\bibitem[35]{published-35} I. Oliveira and R. Santhanam. Conspiracies between Learning Algorithms, Circuit Lower Bounds, and Pseudorandomness. In Proceedings of the 32nd Computational Complexity Conference, CCC’17, Dagstuhl, DEU, 2017. Schloss Dagstuhl–Leibniz-Zentrum fuer Informatik.
\bibitem[36]{published-36} R. Ostrovsky. One-way functions, hard on average problems, and statistical zero-knowledge proofs. In Proceedings of the Sixth Annual Structure in Complexity Theory Conference, pages 133–138, 1991.
\bibitem[37]{published-37} R. Ostrovsky and A. Wigderson. One-way functions are essential for non-trivial zero-knowledge. In Proceedings of the 2nd Israel Symposium on Theory and Computing Systems, ISTCS’93, pages 3–17, June 1993.
\bibitem[38]{published-38} A. A. Razborov and S. Rudich. Natural Proofs. Journal of Computer and System Sciences, 55(1):24–35, 1997.
\bibitem[39]{published-39} H. Ren and R. Santhanam. A Relativization Perspective on Meta-Complexity. In The 39th International Symposium on Theoretical Aspects of Computer Science (STACS 2022), 2022.
\bibitem[40]{published-40} L. Valiant. A Theory of the Learnable. Commun. ACM, 27(11):1134–1142, 1984. doi: 10.1145/1968.1972.
\bibitem[41]{published-41} Thomas Watson. Relativized Worlds without Worst-Case to Average-Case Reductions for NP. ACM Trans. Comput. Theory, 4(3), September 2012.
\bibitem[42]{published-42} H Wee. Finding Pessiland. In Proceedings of the Third Conference on Theory of Cryptography, TCC’06, pages 429–442, Berlin, Heidelberg, 2006. Springer-Verlag.
\bibitem[43]{published-43} D. Xiao. On basing ZK $\neq$ BPP on the hardness of PAC learning. In Proceedings of the 24th Conference on Computational Complexity, CCC’09, pages 304–315, 2009.
\bibitem[44]{published-44} A. Yao. Theory and Application of Trapdoor Functions. In Proceedings of the 23rd Annual Symposium on Foundations of Computer Science, FOCS’82, pages 80–91, November 1982.
\bibitem[45]{published-45} M. Zimand. Efficient Privatization of Random Bits. In IN WORKSHOP ON RANDOMIZED ALGORITHMS, 1997.
\end{thebibliography}
\end{document}
